\part*{Part V --- Price}
\addcontentsline{toc}{part}{Part V --- Price}

\bigskip
\noindent
\emph{A graded clause assigns values to the candidate resolutions of a race. It
does not pick one. Something else has to, and this part is about what that costs.
The distinction between weighing and selecting turns out to be the
quantifier/selection distinction of Escard\'o and Oliva; the condition under which
a selection is definable is a condition on the value algebra; and the strength of
the choice principle needed to resolve a given clause can be written down, one
formula per degree of the Weihrauch lattice, in the same syntactic slot. The part
closes with a conjecture identifying this grading with the packaging grading of
Part II, a failure of subject reduction, and its repair. Source:
\cite{ChoiceFormulae} throughout, with \cite{ChoiceScheduling} for the programme
it answers and \cite{EscardoOliva10} for the categorical backbone.}
\bigskip

\section{Weighing is not selecting}
\label{sec:weighselect}

Write $\Cand(P,r)$ for the set of candidate matches of rule $r$ at term $P$. A
graded clause is a map
\[
  \psi : \Cand(P,r) \longrightarrow \V,
\]
a \emph{weighting} of the fibre. A scheduler is a map
\[
  \varepsilon : \Cand(P,r) \longrightarrow \Cand(P,r) \ \text{(a chosen element)},
\]
a \emph{section}. These are different animals, and conflating them is the error the
whole of this part is organised to avoid.

The categorical statement is due to Escard\'o and Oliva \cite{EscardoOliva10}. A
clause is a \emph{quantifier}: an element of the $K$-monad $(X \to R) \to R$. A
scheduler is a \emph{selection function}: an element of the $J$-monad
$(X \to R) \to X$. There is a canonical morphism $J \to K$ sending
$\varepsilon \mapsto \lambda p.\,p(\varepsilon(p))$, and the statement that a
particular resolver realises a particular clause is the statement that this
morphism sends one to the other.

\begin{principle}[Division of labour]
\label{prin:division}
The clause is the family and its measure. The choice principle is the assertion
that a section exists uniformly. The scheduler is the section. A clause therefore
cannot \emph{be} a choice principle; formulae \emph{certify} an altitude rather
than asserting it.
\end{principle}

This is the reading the programme note asked for. The choice-scheduling
development \cite{ChoiceScheduling} proposes a Curry--Howard correspondence in
which choice principles are types and fair schedulers are terms, and records that
the choice type of a cut is a semantic invariant, not in general computable, and
``approximable from above by types''. The graded clause is that effective
approximation.

\section{Attainment}
\label{sec:attainment}

When is a section definable, so that the resolver can actually hand back a
candidate rather than an expectation?

\begin{proposition}[Attainability]
\label{prop:attain}
\begin{enumerate}[label=(\roman*),leftmargin=2.2em]
\item If $\oplus$ is idempotent and $\psi$ is crisp, a definable section exists.
\item If $\oplus$ is idempotent and $\psi$ is graded, the value is attained
exactly when the winning candidate carries the multiplicative unit.
\item If $\oplus$ is not idempotent, the clause is an \emph{expectation}
quantifier and is never attained into $\Cand$.
\end{enumerate}
\end{proposition}

Measured over $20{,}000$ payoff draws per algebra: $\B$, Viterbi and tropical all
attain $100\%$ of the time on crisp clauses; Viterbi attains $73.83\%$ on graded
clauses; tropical attains $0\%$; $\R_{\ge0}$ attains $0\%$. Clause (ii) was
confirmed at $100\%$ agreement \cite{ChoiceFormulae}.

\begin{remark}[Restriction versus discount]
Gating is free; preference costs. A clause that merely says which candidates are
eligible leaves attainment intact; a clause that says which are \emph{better}
generally does not.
\end{remark}

Non-attainment is not a defect. It is a \emph{lift}: over $\R_{\ge0}$ the
resolution lifts to distributions and is attained in expectation, verified by
Monte Carlo at $0.590826$ against an analytic $0.590000$ and $0.390489$ against
$0.390000$ over $400{,}000$ draws.

\begin{principle}[The lift is the oracle]
\label{prin:lift}
Where a section is not definable, the resolver must be handed something extra ---
a source of randomness, an oracle, an aperture. The type of that extra thing is
what the Weihrauch degree measures.
\end{principle}

Attainment is also where Part II's coalition row was decided. Tropical and Viterbi
are exactly the two algebras where clause (i) holds on the nose; they are also the
two picked out by the quorum semantics; and they are therefore the algebras in
which an observer of a population receives a witness coalition rather than a
scalar (Observation~\ref{obs:threeconstraints}).

\section{The graded fixed point, and three engines of ascent}
\label{sec:engines}

The graded sort as introduced in Part IV is finitary, and a finitary clause cannot
climb. Adding fixed points is what makes altitude available:
\[
  \psi \;::=\; \cdots \;\mid\; X \;\mid\; \mu X.\psi \;\mid\; \nu X.\psi
  \;\mid\; \langle \quot{\ } \rangle_{\V}\, \psi
\]
with the usual guardedness condition. Read $\mu$ as \emph{search} and $\nu$ as
\emph{survival}.

\begin{proposition}[Existence by algebra]
\label{prop:fpexist}
A quantale gives the fixed point by Knaster--Tarski. Over $\R_{\ge0}$ existence
becomes a convergence condition (spectral radius below one). Over $\Cx$ there is
no order and the condition is spectral or, in practice, a budget.
\end{proposition}

Budget is the honest default, and it is a semantic dial: spend budget, gain
altitude.

\begin{proposition}[Three engines]
\label{prop:engines}
There are exactly three ways for a clause to climb: \emph{depth} (nesting a
modality under $\nu$), \emph{width} (quantifying over redex positions under
$\mu$), and \emph{reflection} (the modality $\langle \quot{\ }\rangle$ over quoted
processes). Depth and width are Weihrauch-incomparable --- they correspond to
$\mathsf{C}_{\N}$ and $\mathsf{C}_{2^{\N}} = \mathsf{WKL}$ --- and that
incomparability is the syntactic source of the order being a lattice rather than a
chain.
\end{proposition}

That last clause is worth pausing on, because it converts an assertion into a
derivation. Both \cite{ChoiceScheduling} and the agency note say the order of
choice principles is ``a jumble, not a tower''. Proposition~\ref{prop:engines}
says \emph{why}: there are two independent syntactic engines, and neither
dominates the other.

\section{One formula per degree}
\label{sec:formulae}

The formulae below all live in the same syntactic slot --- a \texttt{where}
clause --- and differ only in shape. This is the concrete sense in which one
programming paradigm covers the whole hierarchy.

\begin{table}[h]
\centering
\small
\renewcommand{\arraystretch}{1.35}
\begin{tabular}{@{}llp{7.1cm}@{}}
\toprule
degree & clause shape & remark \\
\midrule
$\bot$ & fixed-point-free & no choice demanded \\
$\mathsf{LLPO} \equiv \mathsf{C}_{2}$ & two-message race, $\Sigma^{0}_{1}$ guards
  & \rhoc{} programmers write this by accident \\
$\mathsf{C}_{\N}$ & $\mu X.(\lceil\beta\rceil \oplus \langle \Ktx\rangle X)$
  & the width engine \\
$\mathsf{WKL}$ & $\nu X.(\lceil\beta\rceil \otimes \langle \Ktx\rangle X)$
  & the depth engine; compactness \\
$\mathsf{WWKL}$ & the $\R_{\ge0}$ row & Las Vegas; the one point where the two
  axes touch \\
$\lim$ & plasticity acting on the limit
  & learning is free; knowing you have finished is the jump \\
$\mathsf{DC}$ & plasticity itself & no new syntax; an $\omega$-fold product \\
$\mathsf{C}_{\N^{\N}}$ & both engines & \\
\bottomrule
\end{tabular}
\caption{Clause shapes by Weihrauch degree.}
\label{tab:degrees}
\end{table}

Three entries deserve comment.

\paragraph{$\mathsf{LLPO}$ by accident.} A race between two messages whose guards
are semi-decidable is already at $\mathsf{C}_{2}$. This is not an exotic
construction; it is what an ordinary concurrent program looks like. Verified
computationally: over $3 \times 40 \times 60$ instance/stage pairs there were zero
empty supports, and every stage-$k$ selector, for $k = 0..7$ and both tie-breaks,
was defeated by the diagonal that produces at $k{+}1$ on the committed side.

\paragraph{Dependent choice is free.} Proposition~\ref{prop:plasticity} said that a
clause evaluated against the current term makes later values depend on earlier
firings with no extra mechanism. In the present vocabulary that sentence reads:
later fibres depend on earlier selections. It is the principle of dependent
choice, and it required no new syntax at all.

\paragraph{The bottom of the lattice is the extractable fragment.}

\begin{proposition}[Bottom]
\label{prop:bottom}
Every fixed-point-free clause has choice type $\bot$.
\end{proposition}

\begin{corollary}[Two characterisations of one fragment]
\label{cor:extractable}
The fragment a circuit or diagrammatic compiler can consume --- modality-free with
isometric clauses (\S\ref{sec:gradedbuys}) --- and the fragment that demands no
choice coincide. The compiler consumes what it consumes \emph{because} that
fragment demands no choice.
\end{corollary}

\section{Reaching the top, and the two conspicuous routes}
\label{sec:ceiling}

\begin{proposition}[Ceiling]
\label{prop:ceiling}
For a closed term over presented channels with no reflective modality under a
fixed point, the choice type is at most $\mathsf{C}_{\N^{\N}}$.
\end{proposition}

The reason is that names are quoted processes, hence countable, and the store's
bags are finite and presented; a where-clause over a presented store simply cannot
reach full choice. Two routes escape.

\paragraph{Route E: the aperture.} A channel fed by the environment rather than by
the term is not presented. This is the syntactic form of the observation that
capacity lives at the apertures, and it doubles as a security statement: a static
analysis of aperture channels bounds the choice strength of an entire shard.

\paragraph{Route R: reflection.} The clause
$\nu X.(\langle \quot{\ }\rangle_{\V} X \otimes \lceil \beta(\drop{y})\rceil)$ is
impredicative, because the domain of the reflective modality contains processes
that may carry the clause itself.

\begin{corollary}[Trichotomy]
\label{cor:trichotomy}
There are exactly these two routes above the ceiling, and both are syntactically
conspicuous.
\end{corollary}

\begin{remark}[Krivine's \texttt{quote}]
\label{rem:krivine}
Krivine realises dependent choice by adding a \texttt{quote} instruction to the
machine. The \rhoc{} calculus \emph{has} quote as a term former, and has had it
since 2005 \cite{MeredithRadestock05a}. Reflection is the calculus's native choice
primitive. This is the single most compact statement of why the calculus is the
right host for this analysis, and it is also a warning: the same primitive is the
suspect in the failure of subject reduction below.
\end{remark}

\section{Interference is transverse}
\label{sec:transverse}

\begin{proposition}[Transversality]
\label{prop:transverse}
A non-empty interference deficit is not the image of any section, because a
section always lands in the fibre. The complex case is therefore \emph{off} the
Weihrauch lattice, not high on it.
\end{proposition}

Recomputed: phases $0$, $\pi$ and $\pi/2$ give squared coefficients $2.000000$,
$0.000000$ and $1.000000$, while over $\R_{\ge0}$ the minimum over
$100 \times 100$ strictly positive pairs is $0.020000$ and never zero.

Combined with item 5 of \S\ref{sec:negation}, this gives a design rule rather than
a paradox: an aperture may host a resolver of high choice strength, and the
interior may host a coherent resolver, but no single resolver can do both jobs.

\section{Two gradings, one measurement?}
\label{sec:twogradings}

Two Weihrauch gradings have now appeared in this document, derived independently
and for different reasons.

The first is in Part II, from \cite{Zeus}: the degree of $\Pack_{\Coll}$, the
operation packaging a container's splittings back into an object of the same
shape. Its choice residue for $\Mset$ is a linear order on atoms; for $\Pfin$ it
adds decidable equality; for $\List$ it is the identity, because term positions
are already ordered.

The second is this part's: the degree $\chi(c)$ of the choice principle needed to
resolve a race under a given clause. Its ceiling for presented channels is
$\mathsf{C}_{\N^{\N}}$, because names are quoted processes.

\begin{conjecture}[The two gradings agree]
\label{conj:twogradings}
They are the same measurement. The candidate set of a where-clause is a bag ---
the store's produces and consumes --- and selecting a match is packaging that bag
into a distinguished element, i.e.\ $\pack_{\Mset}$. Accordingly the choice type of
a clause is the packaging degree of its candidate container, and
Proposition~\ref{prop:ceiling} is Principle~(trade) of \cite{Zeus}: a presented
store bounds the resolution because term positions are already ordered.
\end{conjecture}

We state this as a conjecture and do not lean on it. The evidence is that the two
bounds have the same content in different words, that both turn on presentedness,
and that both locate their escape in the same place --- a domain that is not
presented, which on one side is an aperture and on the other is an atom object too
large to order syntactically. The obligation is to make ``the candidate container''
precise for clauses using the graded modality, where the candidates at one site
depend on transitions available at another.

\begin{remark}[A related conjecture, better supported]
\label{rem:onethree}
A second identification, verified computationally though not proved, is that the
Escard\'o--Oliva product of selection functions, Weihrauch composition, and
chaining graded sites through the graded modality are one operation. Over
$20{,}000$ tables there were zero disagreements, with a negative control ---
chaining in the wrong algebra --- producing $4{,}483$. The content of the
agreement is that clause algebra and scheduler discipline must match.
\end{remark}

\section{Subject reduction fails, cheaply, and how to repair it}
\label{sec:sr}

\begin{proposition}[Strength can be mailed in]
\label{prop:srfail}
Let $P = c!(Q) \mid \texttt{for}(y \leftarrow c\ \texttt{where}\
\lceil \mathrm{true}\rceil)\,\drop{y}$. Every clause in $P$ is fixed-point-free,
hence at $\bot$ by Proposition~\ref{prop:bottom}; but $P$ reduces to $Q$, which may
be at $\mathsf{WKL}$. Choice type is not preserved by reduction.
\end{proposition}

The mechanism is reflection: substitution replaces $y$ by a quoted process and
$\drop{y}$ drops an arbitrary received process into term position, so a
low-altitude clause can receive a high-altitude one.

\begin{definition}[Stratified \texttt{where}]
\label{def:stratified}
Declare a ceiling $L$ and require a level check on drops: a received process may
be dropped only if its clauses are at or below $L$.
\end{definition}

\begin{proposition}[Subject reduction, stratified]
\label{prop:srfix}
Under Definition~\ref{def:stratified}, choice type is preserved.
\end{proposition}

\begin{corollary}
In a stratified shard, full choice occurs exactly at declared apertures, and the
interior's choice type is invariant.
\end{corollary}

\begin{remark}[The same suspect, four times]
Reflection is now behind: the possibility of length-asymmetric reconvergence and
hence of a second, non-local source of interference; Route R to full choice; the
failure of subject reduction above; and the corresponding question for the graded
coalition modality, where a received process can mail in a high-altitude coalition
formula. Krivine's \texttt{quote} is the same primitive doing the same two jobs ---
supplying choice strength and breaking invariance --- and it is not a coincidence
that a calculus powerful enough to be its own metalanguage has to pay for it here.
\end{remark}

\section{Fairness is a formula}
\label{sec:fairness}

One last consequence, which moves a familiar concept from the typing side of the
correspondence to the term side.

\begin{proposition}[Zero is absorbing]
\label{prop:zeroabsorb}
A continuously enabled receipt whose clause evaluates to zero never fires, and
therefore never acquires the evidence that would make its clause non-zero.
\end{proposition}

This is not hypothetical. In the foraging study, a learner beginning with no prior
evidence has confidence zero, hence clause value zero, hence a receipt that is
permanently enabled and never fires: zero openings, terminal wealth $520.0$, which
is the do-nothing baseline, against $3543.7$ for the same learner given a single
unit of prior evidence \cite{GradedWhere}. Not death, but stasis --- manufactured
by the value algebra.

\begin{definition}[Exploration combinator]
\label{def:explore}
$\psi^{\varepsilon} \;:=\; \psi \oplus \varepsilon$ for a constant
$\varepsilon > 0$: \emph{exploration is a disjunct}.
\end{definition}

\begin{proposition}[Weak fairness is a clause combinator]
\label{prop:weakfair}
$\psi^{\varepsilon}$ enforces weak fairness at the site.
\end{proposition}

Measured: with no disjunct, $0$ firings and terminal $0.00$; with
$\psi \oplus 0.02$, $179$ firings and terminal $83.80$; with one unit of prior
evidence instead, $197$ and $89.20$.

\begin{remark}[Strong fairness has a degree, and the degree is compactness]
Strong fairness is a $\nu$ over the run tree, hence the depth engine, hence
$\mathsf{WKL}$. The classical weak/strong fairness gap thereby acquires a
Weihrauch degree, and the content of the gap is compactness. Over $\Cx$ the
$\varepsilon$-disjunct can cancel, so the fairness repair is a de-tuning device in
the coherent regime and not a guarantee.
\end{remark}

\begin{slogan}
A shard can have any two of: unrestricted reflection, invariant choice types, and
coherent resolution.
\end{slogan}
